
\sourcenote{Genre: core rule (tier 1). Setting as in the \rulename{Ax}
section: valuations, admissibility, entrywise atom satisfaction,
semantic validity; counterexamples read over the one-wire hypernet.}

\textbf{Claim (\texttt{sound\_weaken}).}\enspace The payload supplies a
premise sequent $\Phi \vdash \Psi$ with its sub-derivation, the added
conjunct $\chi$ (\texttt{extra}), and two tree identities: the
conclusion's premise is $\Phi \wedge \chi$ (\texttt{preEq}) and its
conclusion is $\Psi$ (\texttt{postEq}). If $\Phi \vdash \Psi$ is
semantically valid, so is $\Phi \wedge \chi \vdash \Psi$.

\begin{bthm}[soundness of \rulename{Weaken}]\label{thm:weaken-sound}
Under \texttt{preEq} and \texttt{postEq}, semantic validity of the
premise sequent implies semantic validity of the conclusion sequent,
over every hypernet.
\end{bthm}
\begin{proof}
Let $\sigma$ be admissible with $\sigma \models \Phi \wedge \chi$.

\emph{Step 1 (project).} The $\wedge$-clause of satisfaction is
meta-conjunction, so $\sigma \models \Phi$ and $\sigma \models \chi$;
keep the left component and discard the right --- this rule's whole
activity is licensed forgetting.

\emph{Step 2 (apply the premise).} Premise validity at $\sigma$ gives
$\sigma \models \Psi$, which \texttt{postEq} pins as the conclusion's
post.
\end{proof}

\begin{bremark}[\texttt{preEq} forces extension, not replacement]
\label{rem:weaken-preeq}
The identity $\mathrm{seq.pre} = \Phi \wedge \chi$ makes the conclusion's
premise contain the sub-derivation's premise \emph{verbatim} as the left
conjunct of one new $\wedge$-node, with the named $\chi$ on the right.
Drop it and ``weakening'' could \emph{replace} the assumptions instead of
extending them: from $w \leq 0 \vdash w \leq 0$ (an \rulename{Ax}
instance) conclude $\top \vdash w \leq 0$ --- false at $\sigma(w) = 1$.
This context-swap family is the standard counterexample to every dropped
\texttt{preEq} in the structural sections; we cite it from here on
rather than restating it.
\end{bremark}

\begin{bremark}[positional discipline]
$\chi$ arrives as the right child of a single new $\wedge$-node:
$\Phi \wedge \chi$, never $\chi \wedge \Phi$, and never merged into
$\Phi$'s own tree. Any further placement is $\wedge$E/$\wedge$I work.
Naming $\chi$ in the payload keeps the checker's comparison pure
structure --- one tree identity, no matching or search.
\end{bremark}

\begin{bremark}[$\chi$ is unconstrained, harmlessly]
Like \rulename{Ax}, the rule carries no typing side conditions on what
it touches: $\chi$ may be unsatisfiable ($\bot$), ill-typed, or mention
untyped wires. At worst the conclusion's premise becomes harder or
impossible to satisfy --- vacuous validity, never unsoundness.
Weakening only discards information; rules that must \emph{evaluate}
are where typing fields live.
\end{bremark}

\begin{bremark}[monotonicity, locally and wholesale]
In the lattice reading (the $\vee$E section), meet is a lower bound:
the interpretation of $\Phi \wedge \chi$ lies below that of $\Phi$,
which lies below that of $\Psi$. \rulename{Weaken}
is the one-step premise-monotonicity of the consequence relation; the
composition rule's conjunct-wise premise-containment check
(\rulename{Syzygy}'s \texttt{preContains}) is the same monotonicity used
wholesale.
\end{bremark}

\begin{anchors}
Steps 1--2 & \texttt{sound\_weaken}; payload \texttt{PWeakenData} (fields \texttt{extra}, \texttt{preEq}, \texttt{postEq}); wired as \texttt{RuleData.weaken}\\
$\wedge$-clause of satisfaction & \texttt{Formula.holds} (case \texttt{and})\\
lattice reading & \texttt{M3.weaken\_kit}\\
\end{anchors}
